\chapter{质量因子族：重算显著性与有效样本}\label{app:quality}

\section{动机}

第~\ref{ch:result} 章有一个不起眼的结论：本地的质量因子（ROE、毛利率、
现金含量、应计）的旧年度汇总 ICIR 较低。该结果不足以判断全样本显著性。

当时给出的解释是"字段太少"——本地库只有 9 个财务字段。外部快照带来 121 个指标，
其中本地确实拿不到的有：研发费用、商誉、无形资产、有息负债、营运资本、
费用明细、扣非利润、投资收益。

于是有了这一章的实验：\textbf{补充字段是否带来增量预测关联？}

\section{二十个因子}

从 121 个指标里构造 20 个质量因子，统一方向为"数值越大越好"：

\begin{center}
\small
\begin{tabular}{@{}lll@{}}
\toprule
类别 & 因子 & 构造 \\
\midrule
盈利能力 & \texttt{q\_wroe} / \texttt{q\_roe} / \texttt{q\_roa} & 加权 ROE / ROE / ROA \\
 & \texttt{q\_gm} / \texttt{q\_opm} & 毛利率 / 营业利润率 \\
现金质量 & \texttt{q\_cfo\_ni} & 经营现金流 TTM / $|$归母净利 TTM$|$ \\
 & \texttt{q\_accrual} & $-$（归母净利 $-$ 经营现金流）/ 总资产 \\
 & \texttt{q\_cash\_rev} & 经营流入小计 / 营业总收入 \\
研发 & \texttt{q\_rd} & 研发费用 TTM / 营业总收入 TTM \\
商誉 & \texttt{q\_gw} & $-$商誉 / 归母权益 \\
 & \texttt{q\_intang} & $-$（商誉 $+$ 无形）/ 总资产 \\
杠杆 & \texttt{q\_lev} & $-$总负债 / 总资产 \\
 & \texttt{q\_debt} & $-$（短期借款 $+$ 长期借款）/ 总资产 \\
偿债 & \texttt{q\_current} / \texttt{q\_quick} & 流动比率 / 速动比率 \\
营运 & \texttt{q\_wc} & （流动资产 $-$ 流动负债）/ 总资产 \\
 & \texttt{q\_at} & 营业总收入 TTM / 总资产 \\
费用 & \texttt{q\_exp} & $-$（管理 $+$ 销售 $+$ 财务费用）/ 营业总收入 \\
利润含金量 & \texttt{q\_cut} & 扣非净利 / $|$归母净利$|$ \\
 & \texttt{q\_invinc} & $-$投资收益 / $|$归母净利$|$ \\
\bottomrule
\end{tabular}
\end{center}

\section{重算结果}
本节使用补充库、2020--2026 的次日开盘至第六日开盘标签。
调整后 IC 控制对数市值、二十日换手、波动和五日反转。
{\small
\begin{longtable}{@{}lrrrrrr@{}}
\toprule
因子 & 原始 IC & 调整后 IC & 原始 $t$ & 调整后 $t$ & $n$ & $n_{\rm eff}$ \\
\midrule
\endhead
q\_exp & 0.0085 & 0.0185 & 1.70 & 3.97 & 1609 & 410 \\
q\_wroe & 0.0029 & 0.0100 & 0.44 & 2.16 & 1609 & 378 \\
q\_roa & 0.0031 & 0.0083 & 0.51 & 1.68 & 1609 & 376 \\
q\_cut & 0.0064 & 0.0079 & 1.24 & 1.84 & 1609 & 382 \\
q\_accrual & 0.0063 & 0.0079 & 1.93 & 2.56 & 1609 & 363 \\
q\_roe & 0.0033 & 0.0072 & 0.51 & 1.57 & 1609 & 378 \\
q\_opm & 0.0049 & 0.0068 & 0.90 & 1.26 & 1609 & 376 \\
q\_invinc & -0.0047 & 0.0061 & -1.60 & 1.09 & 1609 & 401 \\
q\_intang & 0.0035 & 0.0040 & 1.15 & 1.84 & 1609 & 416 \\
q\_at & 0.0031 & 0.0028 & 0.99 & 0.88 & 1609 & 394 \\
q\_debt & -0.0002 & 0.0017 & -0.06 & 0.65 & 1609 & 425 \\
q\_cfo\_ni & 0.0109 & 0.0010 & 3.27 & 0.24 & 1609 & 371 \\
q\_quick & -0.0000 & 0.0006 & -0.01 & 0.14 & 1609 & 399 \\
q\_rd & -0.0034 & 0.0005 & -0.37 & 0.08 & 1609 & 376 \\
q\_current & -0.0004 & 0.0000 & -0.08 & 0.01 & 1609 & 396 \\
q\_wc & -0.0007 & -0.0000 & -0.12 & -0.00 & 1609 & 389 \\
q\_lev & 0.0012 & -0.0002 & 0.24 & -0.05 & 1609 & 403 \\
q\_gw & 0.0098 & -0.0006 & 2.39 & -0.19 & 1609 & 391 \\
q\_gm & -0.0049 & -0.0038 & -1.00 & -0.82 & 1609 & 415 \\
q\_cash\_rev & -0.0028 & -0.0055 & -0.63 & -1.20 & 1609 & 370 \\
\bottomrule
\end{longtable}
}

原始 IC 与调整后 IC 均由完整日频 IC 序列汇总；$h=5$ 的日频标签重叠，
均值检验使用 Bartlett Newey--West 协方差，滞后阶数为
$L=\max(4,\lfloor4(n/100)^{2/9}\rfloor)$。$n$ 是有值的 IC 日期数，
$n_{\rm eff}=n\gamma_0/\widehat\Omega$ 是原始 IC 均值方差等价的样本量，
不是独立财报次数，也不是未经假设可确定的真实样本数。
完整表另含调整后有效样本量；不同带宽及每五日采样的敏感性见配套 CSV。
本表不再平均各年的 $t$ 值。

\section{撤回“有效样本仅 14 个”的推导}
旧程序先逐年算 $t$，然后求年度 $t$ 的平均。这个平均值不是全样本的检验统计量，
把它与全样本朴素 $t$ 代入 $n(t_{NW}/t_{naive})^2$，得不到可解释的有效样本量。
旧版“20 个因子全部不显著”“预测能力为零”和“质量只适合硬筛选”均不再作为本实验结论。

以 \texttt{q\_cfo\_ni} 为例，本次保存了 \AuditQCfoCount{} 个实际 IC 观测；
朴素 $t=\AuditQCfoNaive$，从同一序列重算的 HAC $t=\AuditQCfoT$，
方差等价样本量约为 \AuditQCfoN{}。这不是“独立季报数”。
慢变因子值与每日收益的横截面相关可以每日变化，不能用财报更新次数替代 IC 相关性测量。
主库的 size、bp、cfp 还包含随股价变动的分母，不能说“一年只变化几次”。

\section{还需要什么证据}
\begin{enumerate}
  \item 对照 5、20、60 日 HAC 带宽及每五日采样结果，检查显著性是否依赖选择；完整结果在敏感性 CSV。
  \item 对全部被尝试的因子和参数进行多重检验；这里的逐因子 $t$ 不是全搜索校正后的显著性。
  \item 核实 PIT 字段、股票池缺失机制和实际收益跨度，再用新的时间段验证。
  \item 质量筛选、长周期持有与主模型增量贡献各自需要实验；本表不证明哪一种用法一定更好。
\end{enumerate}
不显著不能证明无效，显著也不能证明可交易。首先保证统计量来自同一组观测，再解释经济含义。
